\documentclass[11pt]{article}
\usepackage{siunitx}
\usepackage{amsmath}
\usepackage{amssymb}
\usepackage{amsfonts}
\usepackage{subfigure}
\usepackage{amsmath}
\usepackage[x11names]{xcolor}
\usepackage[left=3.5cm,right=3.5cm]{geometry}
\usepackage{framed}
\usepackage{graphicx}
\usepackage{hyperref}
\usepackage{amsthm}
\usepackage{times}
\usepackage{multicol}
\usepackage{booktabs}
\usepackage{tabularx}
\usepackage{multirow}
\usepackage[para,online,flushleft]{threeparttable}

\newtheorem{remark}{Remark}
\colorlet{shadecolor}{Thistle3!50}

\newcommand{\improve}[1]{\textcolor{blue}{#1}}
\newcommand{\highlight}[1]{\textcolor{pink}{#1}}
\newcommand{\todo}[1]{\textcolor{red}{Todo: #1}}
\newcommand{\change}[1]{\textcolor{blue}{#1}}


\def\sss{\scriptscriptstyle{}}
\def\bsss{\sss{\mathrm{B}}}

%%%%%%%%%%%%%%%%%%%%%%%%%%%%%%%%%%%%%%%%

\usepackage[T1]{fontenc}
\setlength{\emergencystretch}{1.5em}
\hypersetup{colorlinks=true,linkcolor=blue,urlcolor=blue,
 pdftitle={Working author response: Multi-robot Coordination under Complex Temporal Logic
Tasks}}
\begin{document}
\title{Authors' Response for the Paper\\
``Multi-robot Coordination under\\Complex Temporal Logic Tasks''}
\author{Meng Guo and Zhongkui Li}
\date{}
\maketitle


\section*{To Associate Editor}
\begin{shaded*}
\noindent\textbf{Q:} \emph{The paper has been well received by the reviewers.
We would kindly ask you to revise the paper according to the reviewers' suggestions
and to resubmit it, \textbf{preferably within December.}
It would be very useful for us if you could highlight in the paper with a different
color the parts that you modified with respect to the current version, and if you
could let us have a reasonably detailed author response.}
\end{shaded*}

\noindent\textbf{A:} We sincerely thank the Associate Editor for handling our
submission and the reviewers for their constructive suggestions. Substantive changes are
highlighted in \textcolor{blue}{blue} in the paper.
The point-by-point replies give the locations and excerpts. The main revisions
are summarized below.

\begin{itemize}
\setlength{\itemsep}{3pt}
\setlength{\parsep}{0pt}
\item \textbf{Motion abstractions and execution.} We unified the weighted transition
system notation, explained how controllers implement symbolic plans, and clarified
the scope of the complexity analysis.

\item \textbf{Temporal logic and planning methods.} We introduced syntactically
co-safe LTL, clarified the link between formulas and automata, and explained
the use of trajectory optimization for bounded-time LTL tasks.

\item \textbf{Coordination algorithms and guarantees.} We added comments to the
algorithms, formalized the knowledge-transfer problem, illustrated the assumptions
on collaboration, and corrected the sufficient penalty bound for task relaxation.

\item \textbf{Organization and examples.} Chapter titles now identify local and
global task specifications, and a recurring service mission connects Chapters~2--6.

\item \textbf{References.} We discussed alternative planning methods and corrected
bibliographic records and citations.
\end{itemize}

\newpage
\section*{To Reviewer 1}

\begin{shaded*}
\noindent\textbf{Q:} \emph{\textbf{Overall Assessment.} The monograph
presents an overview of multi-robot coordination under complex temporal logic tasks. It
comprises seven chapters, including the introduction and ``Outlook and Open Problems.''
Beginning with the fundamentals of temporal logic-based planning for a single robot, it
progressively covers task relaxation to accommodate incomplete or partial knowledge, task
decomposition based on partial orders to enable a group of robots to accomplish a common
mission, and planning under communication constraints. It then considers Markov decision
process (MDP) models of robot motion to address probabilistic task satisfaction. Each
chapter concludes with multiple case studies. Overall, the monograph provides a
comprehensive overview of multi-robot planning and coordination under temporal logic
tasks. The exposition is rigorous and, to the best of my understanding, technically
correct. The following suggestions primarily concern readability, consistency, and the
organization of the material.}
\end{shaded*}
\noindent\textbf{A:} We sincerely thank you for your careful review and constructive
suggestions. We appreciate your assessment of the breadth and technical rigor of the
paper. Your comments on readability, consistency, and organization helped us identify
where clearer notation and stronger connections between chapters would benefit the
reader. Below, we provide point-by-point responses to the comments and a summary of the
changes made in the revision.

\begin{shaded*}
\noindent\textbf{Q:} \emph{\textbf{Suggestions for Improvement.} 1. Reduce the notational
burden and improve
consistency. The presentation is mathematically rigorous, with precise and mostly
unambiguous definitions. However, this rigor comes with a substantial notational burden
that hinders readability, particularly in the later chapters. I found the exposition
especially clear through most of Section 3.1, although Algorithm 3.4 in that section is
already a bit difficult to parse. The notation becomes considerably denser from Section
3.2 onward. Brief comments explaining the purpose of major algorithmic steps would help
readers follow the algorithms (you may consider using the \texttt{\textbackslash comment}
feature provided by many algorithm packages). Changes in notation further compound this
difficulty. Weighted transition systems (WTSs), for example, appear in several chapters,
but their definitions and notation vary. The text immediately following equation (3.20)
states that the WTS structure follows Definition 2.3, yet the two structures are not
identical. Even within the same chapter, the WTS in equation (3.20) is a seven-element
tuple, whereas the one in equation (3.25) contains six elements. Both equations (3.20) and
(3.25) describe mobility abstractions, albeit in slightly different contexts; these
differences do not necessarily require changes to the underlying WTS structure. A unified
structure and consistent notation, with any necessary deviations explicitly explained,
would make it easier for readers, especially those new to the subject, to connect the
concepts across chapters.}
\end{shaded*}
\noindent\textbf{A:} Thank you for pointing out where the notation becomes harder
to follow and for suggesting comments within the algorithms.

\textbf{Notation and symbol meanings.} Your comment about the later chapters
highlights the need to make related quantities easy to identify when notation
becomes more specialized. In Section~5.3.1 (p.~97), for example,
$\widehat{\pi}^{\xi}$ identifies the time and location of each robot's
detachment for request $\xi$. The symbol $\Pi^{\xi}$ maps each detached
robot to its role location. Algorithm~5.5 (p.~102) carries both outputs into
execution. The List of Symbols (p.~143) gives both meanings using the same
symbols as Section~5.3.1.

\textbf{Algorithm readability.} We agree that readers should be able to see
the purpose of each block. The \texttt{\textbackslash tcc} comments in
Algorithm~3.4 (p.~40) mark information exchange and verification of the
current run on the updated model, the choice to resynthesize, repair, or
retain that run, and the update after confirmed arrival. The loop carries
the updated model and product forward. In Section~3.2.3, Algorithm~3.5
(p.~46) checks the horizon before issuing a request. Algorithms~3.7 and~3.8
(pp.~48 and~49) treat the requested action states as terminal targets in the
forward detour searches.

\textbf{Common structure for mobility models.} You are right that the two
mobility abstractions can use the same structure. Definition~2.3 in
Section~2.1.3 (p.~17) and Eqs.~(3.20) and~(3.25) (pp.~41 and~43) use the
same six fields for states, transitions, initial states, propositions,
labels, and weights. In Section~3.2.1 (pp.~41--42), the single navigation
label $\sigma_{\mathcal{M}}$ now belongs to the action set of the composed
model. The table in Section~3.2.2 (p.~43) relates each model's symbols to
the six fields. The time-varying models in Sections~3.1.3 and~3.1.4
(pp.~35 and~38) follow the same order.

\textbf{Model-specific details.} In Section~3.2.2 (p.~43), the robot-specific
weight $T_{\mathcal{M}}^i$ gives the travel duration used to schedule collaborations.
The complete functionality model and conditional action map in Section~3.2.1
(p.~42) include the additional information needed for controller selection
and checking action preconditions, respectively. Section~3.2.1 (p.~43)
keeps each selected action in the projected plan so execution can select
its controller. Section~3.2.2 (p.~43) explains the model extensions and
places the idle action in the local-action set.

We have made the following changes in the revision:
\begin{itemize}
\item \textbf{Notation and symbol meanings.} The definitions in
\change{Section~5.3.1 (p.~97)} read
\change{``The detachment schedule is \(\widehat{\pi}^{\xi}\triangleq
\{\pi_r^{\xi}\}_{r\in\mathcal{N}^{\xi}}\),
where \(\mathcal{N}^{\xi}\) is the set of robots detached from
\(\mathcal{G}^{\texttt{w}}\) to realize \(\mathcal{G}^{\xi}\).
The event \(\pi_r^{\xi}\triangleq(r,t_r,p_r)\) gives the time \(t_r\) and
location \(p_r\) at which robot \(r\) leaves the backbone.
The assignment \(\Pi^{\xi}\) maps each detached robot to its role location
in \(\mathcal{G}^{\xi}\).''}
The corresponding entries in the \change{List of Symbols (p.~143)} describe
$\widehat{\pi}^{\xi}$ as
\change{``Detachment schedule giving the time and location at which each assigned robot
leaves the backbone for request $\xi$.''}
and $\Pi^{\xi}$ as
\change{``Assignment mapping detached robots to role locations for request $\xi$.''}
The text after \change{Eq.~(5.34) in Section~5.3.2 (p.~101)} states
\change{``For an admitted request, each assigned robot uses its detachment event from
\(\widehat{\pi}^{\ell}\) and then moves to its role location under
\(\Pi^{\ell}\).''}
The execution handoff in \change{Algorithm~5.5 (p.~102)} reads
\change{``Store \((\widetilde{\pi}^{\ell},\widetilde{\Pi}^{\ell})\)
            for request \(\xi^{\ell}\).''}
\change{``Execute the stored detachment schedules, route each detached robot
      to its stored role location, and update \(\mathcal{G}\).''}

\item \textbf{Algorithm readability.} The overview in \change{Section~3.1.4
(p.~39)} reads
\change{``During execution, the robot exchanges task-relevant information and checks
its run against the updated model. It then resynthesizes the run, repairs it
locally, or retains it. On confirmed completion in the admissible target set,
the robot records the executed step and selects the next goal.''}
\change{``The updated transition system and product automaton are carried forward
to the next iteration.''}
The three comments in \change{Algorithm~3.4 (p.~40)} are
\change{``Exchange information and check the current run on the updated model.''}
\change{``Resynthesize, repair, or retain the current run.''}
\change{``Record confirmed arrival and advance the goal.''}
In \change{Section~3.2.3 (p.~45)}, the request time is summed over the
original plan indices. \change{Algorithm~3.5 (p.~46)} tests $T_{\textup{next}}>H_i$
before advancing and labels this block
\change{``Request the first assistance needed within the lookahead horizon.''}
The text following \change{Eq.~(3.35) on p.~47} specifies the target rule.
\change{``In the forward search, states in $S_c\setminus S_d$ are forbidden, and states
in $S_d$ are terminal targets. The same rule is used for the forward search
in Algorithm~3.8 with the requester and its target action.''}
The corresponding comments in \change{Algorithms~3.7 and~3.8
(pp.~48 and~49)} read, respectively,
\change{``End each forward detour at its first assisting target.''}
\change{``End each forward path at its first collaboration target.''}

\item \textbf{Common structure for mobility models.} The paragraph following
\change{Definition~2.3 in Section~2.1.3 (p.~17)} states
\change{``A weighted transition system for mobility uses the ordered tuple
$(S,\longrightarrow,S_0,AP,L,w)$, where the entries give the states,
transition relation, initial states, propositions, labeling map, and edge weights.
Equation~(2.4) uses the workspace regions as states, so $S=\Pi$.
The mobility models in Chapter~3 use this order for both
fixed and time-varying abstractions. A time superscript on a component
indicates that its value can change as the robot updates the model.''}
The revised \change{Eq.~(3.20) in Section~3.2.1 (p.~41)} is
{\color{blue}
\begin{equation*}
\mathcal{M} \triangleq
(\Pi_{\mathcal{M}}, \longrightarrow_{\mathcal{M}},
 \Pi_{\mathcal{M},0}, \Psi_{\mathcal{M}}, L_{\mathcal{M}}, W_{\mathcal{M}}),
\tag{3.20}
\end{equation*}}
The accompanying text reads
\change{``The label $\sigma_{\mathcal{M}}$ denotes navigation in the composed model.
Within $\mathcal{M}$, a mobility edge is written as
$\pi\longrightarrow_{\mathcal{M}}\pi'$.''}
The text following \change{Eq.~(3.17) in Section~3.1.4 (p.~38)} explains the
time index.
\change{``The superscript $t$
marks the transition relation, labeling map, and weights available at time
$t \geq 0$.''}
The table following \change{Eq.~(3.25) in Section~3.2.2 (p.~43)} is
introduced as follows.
\change{``The mobility abstractions in (3.20) and (3.25) follow the
same component order as the WTS in Definition~2.3.
Each table column gives the corresponding symbol for one component.''}

\item \textbf{Model-specific details.} The definition following
\change{Eq.~(3.24) in Section~3.2.1 (p.~42)} reads
\change{``The state set is
$\Pi_{\mathcal{R}} \triangleq \Pi_{\mathcal{M}} \times \Pi_{\mathcal{B}}$, and
the action set is
$\Sigma_{\mathcal{R}} \triangleq
\{\sigma_{\mathcal{M}}\} \cup \Sigma_{\mathcal{B}}$.''}
The same subsection explains the fixed-valuation case on \change{p.~42}
and the graph used for synthesis on \change{p.~43}.
\change{``Fixing the external valuation yields a weighted transition system with
action labels on its edges (Baier and Katoen, 2008).''}
\change{``The product search of Chapter~2 uses one
minimum-weight edge for each ordered pair of adjacent states of $\mathcal{R}$
and stores that edge's action label. Parallel edges have the same endpoints
and thus induce the same sequence of state labels. Algorithm~2.1
computes an optimal accepting run $R_{\text{opt}}$ for $\mathcal{A}_{\varphi}$
on this product.
Its projection $\tau_{\text{opt}}$ records the source state, selected action,
and target state of each transition, including the edge that closes the suffix.
For $\sigma_{\mathcal{M}}$, the robot applies the navigation controller in
(2.3) and advances after confirming completion in the admissible target set. For an action
in $\Sigma_{\mathcal{B}}$, it executes the corresponding action controller or
the idle behavior for $\sigma_{\mathcal{B},0}$ and advances on completion.''}
The discussion after the table in \change{Section~3.2.2 (p.~43)} states
\change{``For each retained edge, $T_{\mathcal{M}}^i$ is the nominal travel duration
used to schedule collaborations, obtained from the motion model or estimated from
execution data. The weights $w_c$ and $W_{\mathcal{M}}$ can encode a general
nonnegative motion cost. The complete functionality model in
(3.24) records action labels to distinguish navigation and task
controllers. The conditional action map in (3.23) uses external
workspace valuations to determine which actions are enabled.''}
\change{Section~3.2.2 (p.~43)} also specifies the idle action.
\change{``The local-action set $\Sigma_l^i$ includes the idle action
$\sigma_0 \triangleq \texttt{None}$.''}
\end{itemize}

\begin{shaded*}
\noindent\textbf{Q:} \emph{2. Clarify the relationship between continuous
dynamics and discrete planning. The transition from the continuous dynamics in equation
(2.2) to the discrete representation in equation (2.4) is too abrupt. To my recollection,
the monograph does not subsequently return to the continuous dynamics or explain how to
construct the controller in equation (2.3) to execute the discrete trajectories generated
by the planning algorithms. The discussion following Definition 2.2 is largely cursory and
remains at a high level. One option would be to begin with the discrete representation and
discuss the realization of discrete plans through continuous controllers at the end of
Chapter 2 or later in the monograph. Alternatively, the connection could be developed more
fully when the discrete model is first introduced. As written, readers may wonder how the
transitions $\longrightarrow_c$ in equation (2.4), and those in subsequent WTS motion
models, are precomputed. The complexity analysis excludes the construction of these WTSs
from the underlying dynamics. This is a reasonable choice of scope, but it should be made
explicit. If continuous dynamics are not a focus of the monograph, presenting them later
may improve the flow.}
\end{shaded*}
\noindent\textbf{A:} Thank you for pointing out the abrupt transition between the
robot dynamics and the planning model, and for asking how the resulting plans are
executed.

\textbf{Motion model and transition construction.} We agree with your suggestion to
start with the weighted transition system. The nonlinear dynamics enter when we specify
how a local controller realizes an edge. For example, in Section~2.1.2 (p.~17), the
feedback $u=k(y-x)$ steers a single integrator from a source cell toward an interior
point $y$ in an adjacent target cell. The motion stays in the convex union of the cells
and enters the target in finite time. That controller justifies an edge in the nominal
graph before planning. A robot can add edges online by applying the same feasibility
test to its updated workspace information. Travel durations for collaboration scheduling
come from each robot's motion model or execution data.

\textbf{Plan execution and complexity.} To answer your execution question, the
supervisor advances on confirmed completion in the target region's admissible state
set, which may also constrain heading or speed. If new sensing invalidates the active
edge, the robot cancels the command and reports an interruption. The TIAGo example in
Section~2.4.3 (p.~27) illustrates how navigation and supervisory behaviors carry out
region goals and actions. Section~2.3.5 (p.~23) bounds the size of the translated task
automaton. The search-time and memory bounds assume that automaton and the motion graph
are already available.

We have made the following changes in the revision:
\begin{itemize}
\item \textbf{Motion model and transition construction.}
Section~2.1 (p.~15) opens with
\change{``The planner represents robot motion by a weighted transition system whose
states are workspace regions and whose edges are feasible motions between regions.
Each edge is realized by a local navigation controller. The robot dynamics enter
through the construction and execution of these controllers.''}
Section~2.1.2 (pp.~16--17), titled \change{``Continuous realization of transitions,''}
defines the admissible states and the conditions for retaining an edge.
\change{``For each region $\pi$, let $X_\pi\subseteq\{x\mid p(x)\in\pi\}$ be the
nonempty set of admissible states from which every outgoing controller can start.''}
\change{``For the nominal model, candidate region pairs are tested against a library
of admissible local controllers before task planning. An edge is retained only when
its controller reaches the admissible target set in finite time, respects the input
limits, and remains inside the source--target pair from every admissible source state.''}
Section~3.1.3 (p.~36) applies this rule online.
\change{``An online update adds an edge only after identifying a local controller
that satisfies Definition~2.2 under the updated workspace information. The controller
is stored with the edge.''}
Section~3.2.2 (p.~43) explains the scheduling weights.
\change{``For each retained edge, $T_{\mathcal{M}}^i$ is the nominal travel duration
used to schedule collaborations, obtained from the motion model or estimated from
execution data.''}

\item \textbf{Plan execution and complexity.} Section~2.3.4 (p.~23) states,
\change{``The supervisor advances when the active controller completes and the
estimated state is confirmed in $X_{\pi_{k+1}}$. This target set is also the
admissible start set for the next edge, which makes consecutive controllers
composable.''}
The interruption step in Algorithm~3.4 (p.~40) reads
\change{``Cancel the active command and report an execution interruption.''}
Section~2.3.5 (p.~23) states,
\change{``The task-automaton size bounds quantify the effect of translating the
LTL formula. The search-time and memory bounds below take the task automaton
and finite motion abstraction as inputs. The reported product-construction
times likewise concern the symbolic product. Workspace decomposition,
continuous reachability analysis, controller synthesis, localization, and
low-level tracking are outside these bounds.''}
\end{itemize}

\begin{shaded*}
\noindent\textbf{Q:} \emph{3. Better integrate the discussion of signal
temporal logic. The exposition primarily concerns linear temporal logic (LTL) tasks, but
it introduces signal temporal logic (STL) late in the monograph and moves through the
material quickly. The earlier solutions rely on automata-based approaches, whereas the STL
discussion adopts a different methodology without clearly explaining the transition. A
reader new to the subject may struggle to understand how this discussion relates to the
preceding framework. The connection to the organization of Chapter 6 is also unclear. The
chapter, titled ``Dynamic Uncertainty and Constraints,'' begins with an MDP representation
of the robot and its environment. However, the STL discussion does not appear to involve
the same type of uncertainty, and its connection to the ``constraints'' component of the
title is not clearly articulated. Explaining the motivation for introducing STL, the
reason for the change in methodology, and the relationship to the chapter's themes would
make this material feel more coherent with the rest of the monograph.}
\end{shaded*}
\noindent\textbf{A:} Thank you for pointing out that the link to the earlier LTL
chapters and the role of the final section in Chapter~6 needed clarification.

\textbf{From LTL to bounded-time tasks.} We agree that the earlier LTL formulation
provides a natural starting point. The propositions still record Boolean conditions, such
as whether the team is inside a region. A bounded operator specifies when a condition
must hold. For example, $\mathsf{F}_{[0,h]}a$ is true if $a$ holds at a sample within
$h$ seconds. This is a bounded fragment of metric temporal logic on the sampling grid
(Koymans, 1990). The interval can be expanded into ordinary LTL next steps on that grid
(Alsmann and Lange, 2025). The optimizer uses the same sample indices when it scores
waypoint perturbations, including those inside nested time windows.

\textbf{Trajectory optimization and the role of Section~6.3.} As you note, the
uncertainty in Sections~6.1 and~6.2 concerns transition outcomes and target forecasts,
respectively. Section~6.3 addresses deadlines, formations, and connectivity under a
given task and sampled motion model. The space-time construction of Maity and Baras
(2015) shows how to plan bounded-time tasks using LTL and automata. In our setting,
two trajectories may both satisfy the same formula and offer different formation or
connectivity margins. Direct trajectory optimization uses those margins to rank
candidates, following the quantitative semantics of Fainekos and Pappas (2009) and
Donz\'e and Maler (2010). During the search, the Gaussian-process prior supplies
provisional trajectories for robots that have yet to be planned. Final verification
uses the optimized trajectories of the complete team.

We have made the following changes in the revision:
\begin{itemize}
\item \textbf{Timed task language.} Section~6.3 (p.~124) states,
\change{``Bounded-time LTL here denotes the fragment of metric temporal logic with
bounded future operators evaluated on a uniform sampling grid (Koymans, 1990).''}
Section~6.3.1 (p.~125) defines the bounded until operator as follows.
\change{``The formula $\phi_1\mathsf{U}_I\phi_2$ holds at $t$ if $\phi_2$ holds
at some $j\in t\oplus I$ and $\phi_1$ holds at every sample from $t$
through $j-1$.''}
Section~6.3.2 (p.~127) applies the same timing convention to the perturbation mask.
\change{``The perturbation mask uses the same sample indices as the exact semantics.
At a parent sample $t$, a window $I$ contributes the indices $t\oplus I$.
Nested windows use the sample of their own enclosing operator.''}
Section~6.4.2 (p.~135) gives the time grid for both case studies.
\change{``The 48 samples are indexed from $0$ through $T=47$, so the final sample is
at $4.7\,\mathrm{s}$.''}

\item \textbf{Method and scope.} The Chapter~6 introduction (p.~109) states,
\change{``The first two sections address uncertainty in transition outcomes and
target forecasts, respectively. The third section addresses the timing
and geometry of collective motion under a given task and nominal model.''}
Section~6.3 (p.~124) explains the use of quantitative margins.
\change{``Here the optimizer evaluates the same kind of deadline directly on
candidate trajectories and uses the real-valued predicate margins to rank them.''}
Section~6.3.1 (p.~126) explains how the numerical objective handles constants.
\change{``Before optimization, constants and empty reductions are folded into the
expanded formula tree using the exact min and max rules. If the root reduces
to true, the task robustness term contributes zero to every optimization cost.
If it reduces to false, the sampled task is infeasible and the search returns
failure.''}
Section~6.3.3 (p.~128) defines the signal used for partial plans.
\change{``During local optimization, each robot outside this partial plan is
represented by the mean trajectory of its Gaussian-process prior
conditioned on its endpoints. These means are fixed across the probes.''}
Section~6.3.4 (p.~131) and Algorithm~6.7 (p.~132) place exact verification
before the final ranking. The text states,
\change{``The final ranking includes only nodes with a trajectory for every robot
that pass the verification in Section~6.3.5.''}
Section~6.3.5 (p.~131) states the scope of verification.
\change{``The verification certifies Boolean satisfaction on the sampled trajectory
over the specified finite horizon. Continuous-time safety between samples and
robustness to tracking error require additional interpolation and error bounds, or
direct checks of the realized trajectory.''}
\end{itemize}

\begin{shaded*}
\noindent\textbf{Q:} \emph{4. Use recurring case studies to connect the
chapters. The authors do a good job of presenting case studies at the end of each chapter.
However, the examples change from chapter to chapter, giving the impression of a
collection of simulation results from separate papers. It would be helpful to select one
or two recurring examples and revisit them throughout the monograph. These examples could
demonstrate the qualitative and quantitative differences among the solutions developed in
successive chapters and show how additional modeling assumptions or capabilities affect
the results.}
\end{shaded*}
\noindent\textbf{A:} Thank you for suggesting a recurring example. The same rooms,
corridor, and
service order now provide a common starting point. Chapter~2 asks how a
single robot completes the visits. Chapter~3 adds local knowledge updates
and requests for assistance. Chapter~4 starts with a team-level requirement
and asks which robots should perform the services. Chapter~5 adds decisions
about where and when information is exchanged. Chapter~6 illustrates three
extensions through uncertain passage outcomes, moving targets, and timed
formation requirements.

This lets a reader compare what each method assumes and what it produces.
A quantitative comparison would require a common instance, matched inputs,
and data collected under the same evaluation conditions. The existing case
studies use different task scales and assumptions, so their measurements
cannot establish a cross-method ranking.

Section~1.3.7 (pp.~13--14) introduces the example and its chapter-by-chapter
comparison. The case-study openings in Sections~2.4, 3.3, 4.4, 5.4, and~6.4
(pp.~25, 50, 75, 103, and~133) revisit it. The opening and the scope statement
in Section~1.3.7 read as follows.
\change{Consider the rooms and corridor of Section~2.4.1.
Let $A=r_3$, $D=r_5$, and $E=r_6$ be service regions, with a base at $r_1$.
The initial visitation task is
$\psi\triangleq\Diamond(a\wedge\Diamond d)\wedge\Diamond e$,
where $a$, $d$, and $e$ denote visits to $A$, $D$, and $E$.}

\change{This is a recurring explanatory instance, not a new numerical benchmark.
The later case studies retain their original data and task scales.
Completion time, logical satisfaction, communication load, and robustness
can be compared quantitatively only after fixing a common instance and
the assumptions under which each method is applicable.}

\begin{shaded*}
\noindent\textbf{Q:} \emph{5. Broaden the methodological context and
explain the scope. The exposition could more fully acknowledge alternative approaches. For
example, Eric M. Wolff's doctoral dissertation, Control of Dynamical Systems with Temporal
Logic Specifications [R1], includes methods that use mixed-integer linear programming to
address related problems. I understand that the monograph primarily focuses on
automata-based approaches and therefore need not provide a detailed treatment of every
alternative. However, this choice of scope is applied inconsistently. Throughout the
earlier chapters, the authors focus on automata-based methods despite the availability of
alternative approaches. Yet, in the STL discussion, they focus on a non-automata-based
method even though an automatabased approach is available: Maity and Baras [R2] formulate
bounded-time temporal logic specifications using LTL in a space-time representation,
enabling the use of automata-based planning tools. It is therefore unclear why the authors
depart from their chosen methodological focus in the final subsection when an approach
exists that could preserve the coherence of the exposition.}
\end{shaded*}
\noindent\textbf{A:} Thank you for pointing us to Wolff (2014) and Maity and Baras (2015).
Both
references help explain why an automaton is useful in some parts of the paper
and why the final section works directly with trajectories. Chapters~2--5
use accepting runs and their temporal dependencies as shared objects for
plan repair, task decomposition, and coordination. Wolff's mixed-integer
encodings provide a different way to connect logical requirements and
trajectory variables. The relative cost of the two formulations depends on
the dynamics, horizon, and specification.

You are also right that bounded time alone does not dictate the choice of
method. In a space-time formulation, time is represented in the state and the
planner can use LTL tools, as in Maity and Baras. Section~6.3 optimizes
quantitative formation, separation, and connectivity margins on sampled
trajectories. Those quantities motivate the trajectory representation.

Section~1.2.2 (p.~4) discusses Wolff's two formulations. Section~6.3
(p.~124) explains the space-time alternative and why the section uses
trajectory optimization.
\change{Wolff (2014) develops both automaton-guided constrained
reachability and direct mixed-integer encodings for dynamical systems.
The latter encode logical requirements together with trajectory variables
instead of first constructing a full finite motion abstraction.}

\change{Maity and Baras (2015) use LTL and a space-time abstraction to plan
with bounded-time requirements. Here the optimizer evaluates the same kind
of deadline directly on candidate trajectories and uses the real-valued
predicate margins to rank them.}

\begin{shaded*}
\noindent\textbf{Q:} \emph{\textbf{Technical Comments.} Most of the
analysis appears correct to me. The following questions identify points that would benefit
from clarification.}
\end{shaded*}
\noindent\textbf{A:} Thank you for identifying the points in the analysis that need
clarification.
We appreciate the opportunity to explain the assumptions and the penalty
choice more precisely.

\begin{shaded*}
\noindent\textbf{Q:} \emph{1. In Definition 2.2, must the set $\pi_i$ be
connected? More generally, what topological assumptions are required?}
\end{shaded*}
\noindent\textbf{A:} Thank you for asking about the topology of $\pi_i$. Definition~2.2
does
not impose connectedness as a separate condition. It requires a controller
that can reach the admissible target set from every admissible source state while
remaining in the source--target pair. A disconnected source can satisfy
that requirement if every component admits the required confined motion
under the controller. Conversely, even a connected region can fail the
requirement because of the dynamics or input bounds. The abstraction also
needs measurable regions with well-defined membership and uniform labels.
When a component or an arrival orientation prevents a valid transition,
the abstraction needs to retain that information.

The discussion after Definition~2.2 (p.~17) specifies the region assumptions
and explains what to do when the controller only works on part of a region.
\change{The regions are taken to be measurable and to have well-defined membership
tests and uniform labels. Connected, geometrically simple cells are useful
for constructing navigation laws, but Definition~2.2
requires the stronger control property stated there. If a disconnected
region contains an unreachable component, an edge cannot be justified by a
controller that works only from another component. The abstraction must
then split the region or retain the relevant entry-state information.
The admissible state sets may restrict orientation or velocity. If a common
set cannot support the intended outgoing edges, the abstraction must retain
the relevant state information through a finer partition of the state space.}

\begin{shaded*}
\noindent\textbf{Q:} \emph{2. How is the Eval function defined in Section
3.1.1, approximately three lines above equation (3.2)?}
\end{shaded*}
\noindent\textbf{A:} Thank you for pointing out that the encoding was underspecified. The
input
$l$ is one label, a subset of $AP$. For a fixed ordering of the $K$ propositions
in $AP$, $\texttt{Eval}(l)$ is the length-$K$ vector that records membership
in $l$. The $\ell_1$ distance between two such vectors therefore equals the
number of propositions on which the labels disagree.

The definition before Eq.~(3.2), in Section~3.1.1 (p.~30), now reads
\change{$\texttt{Eval}(l)_k\triangleq 1$ if $a_k\in l$ and $0$ otherwise, for
$k\in\{1,\cdots,K\}$.}

\begin{shaded*}
\noindent\textbf{Q:} \emph{3. What motivates the linear combination in
equation (3.3)? Other ways of combining task feasibility and environment transitions seem
possible. A brief explanation of this choice would be helpful.}
\end{shaded*}
\noindent\textbf{A:} Thank you for asking why the two quantities are combined linearly.
The
motion term prices an executable edge, and the specification term measures
how far its observed label is from enabling the intended automaton
transition. The coefficient $\alpha$ prices a unit of that mismatch in
motion-cost units. With nonnegative additive edge weights, the relaxed
product can use the same shortest-path routines as the nominal product.
Other choices are possible. A lexicographic objective gives violation
strict priority, and a proposition-weighted distance can give certain
mismatches greater importance. In Section~3.1.2, hard requirements are
enforced by the transition relation itself.

The paragraph following Eq.~(3.3) in Section~3.1.1 (p.~30) explains the
choice and two alternatives.
\change{The additive form in (3.3) retains nonnegative edge weights,
so the same shortest-path routines can compare relaxed accepting runs.
The two terms have different roles: $w_c$ measures implementation cost,
whereas $\texttt{Dist}$ counts proposition mismatches at an automaton
transition. The coefficient $\alpha$ converts one unit of this distance
into the units of the motion cost. This is a design choice, not a unique
consequence of the task semantics. A lexicographic objective can instead
prioritize violation before motion cost, and other violation measures can
encode proposition priorities.}

\begin{shaded*}
\noindent\textbf{Q:} \emph{4. In Theorem 3.1, the threshold
$\underline{\alpha}$ is unknown. Is an estimate of this threshold needed to select
$\alpha$? If so, how can such an estimate be obtained? Choosing $\alpha$ ``as large as
possible'' may be a practical heuristic, but a more informative guideline would be
valuable. In addition, quantities used in a theorem statement should preferably be defined
before or within the statement, rather than introduced only in the proof, as is the case
for $\underline{\alpha}$.}
\end{shaded*}
\noindent\textbf{A:} Thank you for asking how to select $\alpha$. Any feasible accepting
prefix-suffix run of the standard product gives a sufficient bound. Its
cost $C_{\textup{feas}}$ can be computed with the same suffix weight $\gamma$
used in the relaxed objective. A violation contributes at least $1$ on the
prefix or $\gamma$ on the suffix, so a violating run costs at least
$\alpha\min\{1,\gamma\}$. Choosing $\alpha$ strictly above
$C_{\textup{feas}}/\min\{1,\gamma\}$ makes such a run more expensive than the
feasible witness. Accepting-cycle detection can produce a witness without
requiring an optimal nominal run. If a time-limited search yields no witness,
a parameter sweep can still describe the observed trade-off. This bound requires
a feasible witness. Increasing
$\alpha$ cannot restore satisfaction for an infeasible task.

Theorem~3.1 and Eq.~(3.7) (p.~32) define the threshold in the theorem
statement. The relevant passage reads
\change{Suppose that $\gamma>0$ and that $\varphi$ is feasible over
$\mathcal{T}_c$. Let $R_{\textup{feas}}$ be any accepting prefix-suffix
run of the standard product $\mathcal{A}_p$, let
$C_{\textup{feas}}\triangleq J(R_{\textup{feas}})$ with $J$ defined in
(2.13), and define
\begin{equation*}
\underline{\alpha}\triangleq
\frac{C_{\textup{feas}}}{\min\{1,\gamma\}},
\tag{3.7}
\end{equation*}
where $C_{\textup{feas}}$ is finite and nonnegative.
If $\alpha>\underline{\alpha}$, then the balanced plan
$\tau_{\textup{bal}}$ satisfies $\varphi$.}

Equations (3.5) and (3.6) (p.~31) and Eq.~(3.12) (p.~34) include the suffix closing edge
in their respective sums. The closing-edge
terms are \change{$w_c(\pi_n,\pi_f)$},
\change{$\texttt{Dist}(L_c(\pi_n),\chi(q_n,q_f))$}, and
\change{$\texttt{Dist}(L_c(\pi_n),\chi_2(q_n'|_{Q_2},q_f'|_{Q_2}))$}. The text after
Eq.~(3.6) (p.~31) states
\change{The closing edge of the suffix is included,
as in (2.13). Positivity of $\gamma$ is required for
zero accumulated distance to imply that every recurrent transition
satisfies the original automaton constraint.}

\begin{shaded*}
\noindent\textbf{Q:} \emph{5. Could the authors elaborate on Assumption
3.1? Its role in the analysis is understandable, but it is less clear when the assumption
is satisfied in practice and when it is likely to be violated. Examples of both situations
would help readers assess its applicability.}
\end{shaded*}
\noindent\textbf{A:} Thank you for asking when this assumption is realistic.
Assumption~3.1
requires the assignment problem for a request made at time $t_m$ to become
feasible by $t_m+T$, for a common finite $T$. A dedicated helper pool can
meet that requirement when each role has a reachable helper, each existing
commitment ends within a known bound, communication and travel delays are
bounded, and requests are served fairly. The bound needs to include both
helper release and receipt of feasible replies. Permanent loss of a
required capability or link, circular waiting, and starvation can all
invalidate the assumption. In those cases the completion guarantee in
Theorem~3.8 does not apply.

The discussion following Assumption~3.1 in Section~3.2.3 (p.~49) gives
the concrete conditions and failure examples.
\change{Assumption~3.1 is a bounded-availability condition,
not a consequence of having few collaborations. It holds, for example,
when every requested assisting role has a reachable helper, every helper's
existing commitment has a uniform finite duration, request handling is
fair, and travel and message delays are uniformly bounded. A dedicated
helper pool serving finitely many requests under those conditions is a
representative instance. The bound $T$ must cover both the release of
busy helpers and the time needed to obtain feasible replies.
The assumption fails if a required capability is absent, a blocked route
permanently separates a helper, communication is lost indefinitely, or
robots wait cyclically for one another while rejecting requests.
It can also fail through starvation under an unfair request order.}

\begin{shaded*}
\noindent\textbf{Q:} \textbf{References.}

\noindent [R1] Eric M. Wolff. \emph{Control of Dynamical Systems with Temporal
Logic Specifications}. Ph.D. dissertation, California Institute of
Technology, 2014. doi:10.7907/TGFR-SS39.

\noindent [R2] Dipankar Maity and John S. Baras. ``Motion planning in dynamic
environments with bounded time temporal logic specifications.'' In
\emph{2015 23rd Mediterranean Conference on Control and Automation (MED)}.
IEEE, 2015. Author-hosted paper.
\end{shaded*}
\noindent\textbf{A:} Thank you for these references.
\href{https://thesis.caltech.edu/8078/}{Wolff (2014)} provides the
mixed-integer example discussed in Section~1.2.2 (p.~4), and
\href{https://johnbaras.com/wp-content/uploads/2020/10/15-08-MCC-2015-Final-Pub-Ver.pdf}%
{Maity and Baras (2015)} provide the space-time approach in Section~6.3
(p.~124). The corresponding bibliography records appear on pp.~158 and~154,
respectively.

We would like to thank you again for the time and care you devoted to
reviewing our paper. Your questions and constructive feedback helped us explain
the modeling assumptions, algorithmic choices, and scope of the results more
clearly.

\newpage
\section*{To Reviewer 2}

\begin{shaded*}
\noindent\textbf{Q:} \emph{The article
focuses on multi-robot coordination with temporal logic specification. It starts from a
single robot case, then moves to coordination among robots and how to decompose the
original global task to individual robots. After these fundamentals, it then considers
imperfections such as limited knowledge and communication and uncertainties. Finally, it
concludes with open problems. The structure of the article is suitable and builds things
up very well. What I particularly like is that each section has a well-defined problem
formulation. I have some minor comments below.}
\end{shaded*}
\noindent\textbf{A:} Thank you for the careful assessment and for noting the progression
from
single-robot foundations to coordination with uncertain information. Your
specific comments helped us identify where the problem formulations and
notation needed more context.

\begin{shaded*}
\noindent\textbf{Q:} \emph{Section 2.1.1: I do not see the need to
introduce $AP_r$ the way it is introduced, i.e., with each proposition indicating whether
the robot is in the corresponding region. This unnecessarily leads to exponential
explosion. With n regions, we need $2^n$ propositions to represent it. However, a robot
can only be in 1 region at a time (as correctly noted).}
\end{shaded*}
\noindent\textbf{A:} Thank you for raising the distinction between propositions and
valuations.
With $N$ regions, there are $N$ location propositions, one per region.
Although the Boolean alphabet $2^{AP_r}$ contains $2^N$ subsets, a robot
occupies exactly one region at a symbolic step. Only the $N$ singleton
location valuations can arise. Workspace properties can coexist with that
one location proposition. The region propositions let a formula name
individual destinations directly, and the motion abstraction uses $N$
location states.

The paragraph following the definition of $AP$ in Section~2.1.1 (p.~15) makes this
explicit.
\change{The location encoding uses $N$ propositions, not $2^N$ propositions.
Although the general alphabet $2^{AP_r}$ has $2^N$ subsets, only the
$N$ singleton location valuations are admissible under the partition
assumption. The motion abstraction therefore has $N$ location states;
it need not enumerate the full Boolean alphabet. Workspace-property
propositions may hold together with the unique location proposition.
At a shared geometric boundary, a fixed ownership convention assigns the
boundary point to one region, preserving a unique location valuation.}

\begin{shaded*}
\noindent\textbf{Q:} \emph{Definition 2.2: Please specify whether $t_s$ and
$t_f$ are given or we should interpret this as they exist.}
\end{shaded*}
\noindent\textbf{A:} Thank you for catching the missing quantifiers. The start time $t_s$
is
arbitrary. Given $t_s$ and an admissible initial state, the same edge
controller must yield some finite arrival time $t_f>t_s$. The arrival
time may depend on that execution. The condition applies from every
admissible source state so that symbolic transitions can be composed.

Definition~2.2 (p.~16) specifies the quantifiers with Eq.~(2.3).
\change{For every admissible initial state $x_s\in X_{\pi_i}$ and starting time
$t_s\geq 0$, there exists a finite arrival time $t_f>t_s$ such that, with
$x(t_s)=x_s$, the control input is given by
$u(t)\triangleq\mathcal{U}\bigl(x(t),\pi_i,\pi_j\bigr)$, $t\in[t_s,t_f]$,
where the input respects its limits and the closed-loop trajectory of
(2.2) reaches $x(t_f)\in X_{\pi_j}$, with $p(x(t))\in\pi_i\cup\pi_j$ for all
$t\in[t_s,t_f]$.}

\begin{shaded*}
\noindent\textbf{Q:} \emph{The statement after Definition 2.2: "The latter
prevents a symbolic transition...". It's unclear what "latter" means in this
context.}
\end{shaded*}
\noindent\textbf{A:} Thank you for pointing out the ambiguous reference. It refers to
confinement
for the entire motion segment. Even with a valid controller and an endpoint
in the target, a trajectory that crosses a third region would introduce
an unmodeled visit.

The paragraph immediately after Definition~2.2 (p.~16) names the condition.
\change{The confinement condition
$p(x(t))\in\pi_i\cup\pi_j$ prevents a symbolic transition from implicitly
using intermediate regions absent from the discrete plan.}

\begin{shaded*}
\noindent\textbf{Q:} \emph{The statement after Definition 2.3: "Temporal
logic formulas are evaluated over these infinite words..." I believe the notion of "words"
has not been introduced.}
\end{shaded*}
\noindent\textbf{A:} Thank you for flagging the missing introduction. A word in this
setting is a
sequence of region labels. Each symbol is the set of propositions true at a
single symbolic step, and concatenating those symbols along an infinite path
gives its trace. Definition~2.5 (p.~18) then uses that trace alphabet to state
the LTL semantics.

The paragraph after Definition~2.3 in Section~2.1.3 (p.~17) introduces the
term before defining the trace.
\change{An infinite word over an alphabet is an infinite sequence of symbols from
that alphabet. Here the alphabet is $2^{AP}$, so each symbol is the set of
propositions true at one symbolic step. Every path induces such a word.}

\begin{shaded*}
\noindent\textbf{Q:} \emph{Definition 2.8: The way this definition is
stated makes Lemma 2.1 incorrect. In particular, the way Definition 2.8 is current stated
does not impose any condition that $\mathcal{A}_{\varphi}$ has to correspond to $\varphi$.
If we simply follow the definition, we can define Q, $\delta$, $Q_0$ and $\mathcal{F}$
arbitrarily since the definition does not relate these elements to $\varphi$. But then
Lemma 2.1 states that this arbitrary construction ensures that the accepted language of
$\mathcal{A}_{\varphi}$ is the language generated by $\varphi$.}
\end{shaded*}
\noindent\textbf{A:} Thank you for distinguishing the structural definition from the
language
equivalence claim. You are right that arbitrary choices of $Q$, $\delta$,
$Q_0$, and $\mathcal{F}$ need not define an automaton for $\varphi$.
The equivalence follows from an LTL-to-B\"uchi translation, as in
Theorem~5.37 of Baier and Katoen (2008). Thus $\mathcal{A}_{\varphi}$
in the synthesis arguments denotes a translation result. In the finite
product, a reachable accepting cycle then supplies a prefix-suffix run.

Definition~2.8 and Lemma~2.1 in Section~2.3 (pp.~19--20) state the structural
and language claims separately.
\change{A Nondeterministic B\"uchi automaton over $2^{AP}$ has the tuple structure
below. Its components specify an automaton, but do not by themselves
associate it with any particular LTL formula:}

\change{For every LTL formula $\varphi$ over $AP$, there exists a
Nondeterministic B\"uchi automaton $\mathcal{A}_{\varphi}$ with the structure
of Definition~2.8 such that
$\mathcal{L}_{\omega}(\mathcal{A}_{\varphi})
=\textup{\texttt{Words}}(\varphi)$.}

\begin{shaded*}
\noindent\textbf{Q:} \emph{Sections 3.1.1 and 3.1.2: These sections are
about what to do when the specification $\varphi$ is not satisfiable. It's unclear to me
what this has to do with coordination. I feel like it fits Chapter 2 better.}
\end{shaded*}
\noindent\textbf{A:} Thank you for the suggestion. Sections~3.1.1 and~3.1.2 address
infeasibility
for a single robot. Their role in Chapter~3 is to supply the relaxed cost
and hard-task transition constraints that each robot uses later when
revising its plan after an information exchange. Placing these sections in
Chapter~2 would put that material with the single-robot foundations. The
current organization keeps it immediately before the multi-robot use,
with a reading path for readers already familiar with relaxation.

The opening of Section~3.1 (p.~29) explains the scope and gives that
reading path.
\change{Sections~3.1.1--3.1.3 are single-robot
prerequisites, not coordination algorithms by themselves. They are retained
here to introduce the local update mechanism immediately before its
multi-robot use in Section~3.1.4. Readers familiar
with task relaxation may proceed directly to that subsection.}

\begin{shaded*}
\noindent\textbf{Q:} \emph{Section 3.1.3: Similar to 3.1.1 and 3.1.2, this
section has nothing to do with robot coordination. Consider moving it to Chapter
2.}
\end{shaded*}
\noindent\textbf{A:} Thank you for raising the same issue for online adaptation.
Section~3.1.3
describes how one robot incorporates its observations, updates the local
model, and finds a safe continuation compatible with the executed history.
Section~3.1.4 uses that same process after the robot receives task-relevant
labels from neighbors. This dependency is the reason Section~3.1.3 sits
directly before the collaborative extension in the current organization.

The signpost at the start of Section~3.1 (p.~29) covers Section~3.1.3 as
part of the single-robot material. The text following Problem~3.7 in
Section~3.1.4 (p.~38) specifies what the extension transfers.
\change{The shared information concerns region labels; the sender's
motion feasibility is not transferred to a receiver with different
capabilities.}

\begin{shaded*}
\noindent\textbf{Q:} \emph{Section 3.1.4: Unlike earlier sections, this
section does not have a precise problem formulation.}
\end{shaded*}
\noindent\textbf{A:} Thank you for asking for a precise formulation. Each robot starts
with its
own motion model and hard and soft tasks. Local observations and messages
from current neighbors can change its labels. The coordination problem is
to decide what information to exchange and how to update each robot's plan so
it remains valid and satisfies the hard task relative to the current model.
Exact resynthesis minimizes the local relaxed cost. The guarantee depends
on there being a safe accepting continuation that agrees with what the
robot has already executed. Theorem~3.6 (p.~41) applies during uninterrupted
execution and requires each necessary search or repair to succeed in time.

Problem~3.7 in Section~3.1.4 (p.~38), before the initial synthesis,
states the inputs, decision, objective, and feasibility premise.
\change{Given the local models $\{\mathcal{T}_i^0\}$, local hard--soft tasks
$\{\varphi_i\}$, time-varying neighbor sets $\{\mathcal{N}_i^t\}$,
and local sensing records $\{\textbf{Sense}_i^t\}$ defined below, design
a distributed information-exchange and plan-update scheme. Each robot
must use its own observations and received task-relevant labels to
maintain a valid plan satisfying its hard specification on its current
model, while minimizing the local relaxed cost when exact resynthesis
is invoked. The requirement is conditional on the existence of a safe
accepting continuation compatible with the executed history after each
model update.}

\begin{shaded*}
\noindent\textbf{Q:} \emph{Section 3.3: This example uses sc-LTL while
Chapter 2 only introduces LTL. Would it be possible to include in Chapter 2 the definition
of these LTL fragments that are used later so readers understand their differences? In
fact, the statement after Equation (4.1) says "where the syntax and semantics of sc-LTL
follow Chapter 2" but Chapter 2 only gives LTL. I understand that the syntax of sc-LTL is
exactly the same as LTL as an sc-LTL formula is also an LTL formula but the article does
not make this obvious for those who don't know these logics.}
\end{shaded*}
\noindent\textbf{A:} Thank you for suggesting that the fragment be introduced before the
examples
and for pointing out the cross-reference after Eq.~(4.1). Every syntactically
co-safe LTL formula is an LTL formula with the same infinite-word semantics.
The syntax restricts negation to atoms in positive normal form and uses
next and until, with eventually as shorthand. Its practical benefit is
the finite good prefix: once the prefix is observed, every infinite
continuation satisfies the formula. Ordered visitation can therefore be
certified at a finite step, whereas persistent surveillance cannot. Synthesis
can use a good-prefix automaton or an equivalent accepting-absorbing
construction. An arbitrary B\"uchi accepting state alone is insufficient.

Section~2.2.1 (p.~19) defines the grammar and explains the finite certificate.
\change{The syntactically co-safe fragment, denoted by $\mathrm{sc\mbox{-}LTL}(AP)$,
is generated in positive normal form by
\[
\psi ::= \top\mid\bot\mid a\mid\neg a
\mid\psi_1\wedge\psi_2\mid\psi_1\vee\psi_2
\mid\bigcirc\psi\mid\psi_1\mathsf{U}\psi_2,
\]
where $a\in AP$; $\Diamond\psi$ is shorthand for
$\top\,\mathsf{U}\,\psi$. Negation is restricted to atomic propositions.
Every sc-LTL formula is an LTL formula and retains the infinite-word
semantics of Definition~2.5.
Its distinguishing property is that every satisfying infinite word has a
finite good prefix: a finite word $u$ such that $uv\models\psi$ for
every infinite continuation $v$. The notation $u\models\psi$ for a finite
execution in later chapters means that $u$ is such a good prefix; it does
not introduce a separate finite-trace logic.}

The text following Eq.~(4.1) (p.~56) points to that definition.

\change{where $\mathrm{sc\mbox{-}LTL}(AP)$ is the syntactically co-safe fragment
defined in Section~2.2.1, here with
$AP=\mathbf{p}\cup\mathbf{a}\cup\mathbf{c}$.}

Section~3.2.2 (p.~44) uses the term \change{sc-LTL formula} for the
dependent local task.

\begin{shaded*}
\noindent\textbf{Q:} \emph{Chapter 3 vs Chapter 4 naming: The titles of
Chapter 3 and Chapter 4 seem to imply that they're solving a similar problem (one using
bottom up approach and the other using a top-down approach). But they're solving very
different problems. Chapter 3 does not care about global specifications at all. It's just
about how the robots coordinate so that each of them accomplishes their individual task
$\varphi_i$ or a relaxed version of it but this coordination and relaxation does not
really take into account the global specification. I'd just avoid using "bottom-up" and
"top-down" in the title.}
\end{shaded*}
\noindent\textbf{A:} Thank you for pointing out the distinction between the two input
problems.
Chapter~3 begins with prescribed local formulas $\varphi_i$. Coordination
helps each robot fulfill its own requirements, for example by sharing
observations or assigning a helper for a collaborative action. Chapter~4
begins with one team-level formula. Task decomposition and assignment must
then preserve the temporal dependencies in that formula. The direction
of information flow is secondary to this difference in the specifications
supplied to each method.

The revised titles identify those inputs directly on pp.~29 and~55.
Chapter~3. \change{Coordination under Local Temporal Logic Tasks}

Chapter~4. \change{Synthesis and Coordination under Global Temporal Logic Tasks}

The opening of Chapter~3 (p.~29) also states the scope of its guarantees.
\change{No global formula is decomposed in this chapter. The terms local and global
refer to where specifications are supplied, not to two interchangeable
algorithms for the same input. The guarantees concern the supplied local
tasks and their stated interaction assumptions; they do not certify an
additional, unspecified team mission.}

\begin{shaded*}
\noindent\textbf{Q:} \emph{Section 4.1.1: The notation of Eq (4.1) has not
been introduced. On the other hand, Definition 2.4 introduces an LTL formula over AP. You
can introduce the notation of Eq (4.1) here, i.e., say that this is ltl(AP). Or you can
just remove this notation and just say in Equation 4.1 that you consider sc-LTL formula
over $\mathbf{p}\cup\mathbf{a}\cup\mathbf{c}$.}
\end{shaded*}
\noindent\textbf{A:} Thank you for pointing out that the formula-set notation was used
before
it was defined. $\mathrm{LTL}(AP)$ denotes the set of LTL formulas
over $AP$, and $\mathrm{sc\mbox{-}LTL}(AP)$ denotes its syntactically co-safe
fragment. In Eq.~(4.1), the alphabet comprises the location, local-action, and
collaborative-behavior propositions defined for the global task. Writing the formula
as an element of this set makes both the fragment and the alphabet explicit.

Section~2.2.1 (p.~19) introduces the formula set, and the text immediately
after Eq.~(4.1) in Section~4.1.1 (p.~56) specializes it to the mission
alphabet.
\change{The set of LTL formulas over $AP$ is denoted by $\mathrm{LTL}(AP)$.}

\change{where $\mathrm{sc\mbox{-}LTL}(AP)$ is the syntactically co-safe fragment
defined in Section~2.2.1, here with
$AP=\mathbf{p}\cup\mathbf{a}\cup\mathbf{c}$.}

We sincerely thank you again for your time, effort, and careful reading
of our paper. Your comments helped us clarify how the definitions, problem
formulations, and chapter organization fit together. We appreciate your
constructive guidance throughout the review.

\newpage\section*{Bibliographic verification --- supplemental author note}
This additional verification follows the authors' request to audit the
references. It is not an additional reviewer comment.

All 528 original bibliography records were checked, including 233 originally
cited records. The revised paper cites 230 verified records. Seven
unverified cited records have been replaced with sources suited to the
surrounding claims; four new journal papers are introduced and the other
replacements reuse verified existing works. The new works appear in
\emph{Science Robotics}, \emph{The International Journal of Robotics
Research}, and \emph{Proceedings of the IEEE}. Authentic books, original
conference papers, and software references remain where appropriate.

The active database contains 450 verified records. A separate review
bibliography preserves all 82 held original records unchanged: seven now
replaced cited records, 74 unresolved unused records, and one confirmed
retracted unused paper. The 74 unresolved records are not declared
fabricated. Each has an explicit author question in
\texttt{REFERENCE\_AUTHOR\_QUESTIONS.md}. The complete evidence and
field-level corrections are supplied in \texttt{REFERENCE\_AUDIT.md} and
\texttt{REFERENCE\_AUDIT.json}.

Notable corrections include the STRIPS publication year, the journal
locator for \emph{Action Languages}, the actual IEEE publication of the
construction-robotics paper, Swamy et al.'s actual \emph{Scaled Autonomy}
paper, publication types, missing authors, final pagination, and registered
DOIs. The MoRoCo citation is pinned to arXiv version~1, matching the title
and authors used in the supplied chapter. Changed bibliography entries,
citations with changed author/year, replacement citations, and revised
prose are blue.

The Retraction Watch screening covered a dated dataset of 73,477 records.
The unused whale-optimization paper is confirmed retracted and excluded
from scientific support. A title-only hit for ``Dynamic Programming'' was
an unrelated Alan Sultan chapter, not Bellman's Science paper. No revised
cited work matched a confirmed retraction in this screening. This result
does not certify unindexed or future notices.

The source-package audit accounts for every original project file and
bibliography record. All original figure assets and labels are retained.
The final compilation and integrity results are in
\texttt{VALIDATION.md} and \texttt{SOURCE\_INTEGRITY.md}.

\newpage\section*{Author questions --- remove before submission}
\noindent These notes are for the authors. They identify decisions and technical passages
deliberately left unchanged. They are not claims that the corresponding issues have been
resolved.
\par\medskip
\noindent\textbf{AQ1 --- Organization of single-robot extensions}\par
\todo{Should Sections~3.1.1--3.1.3 move into Chapter~2? The present revision
adds a prerequisite signpost but leaves their order unchanged. This decision
affects Reviewer~2 comments~6 and~7.}
\par\medskip
\noindent\textbf{AQ2 --- Quantitative recurring benchmark}\par
\todo{Should the recurring nine-region service example become an experimental
benchmark across the applicable methods? Which code, common robot model,
task instances, and logs should be used? The added example is explanatory;
there are no fabricated measurements or cross-method rankings.}
\par\medskip
\noindent\textbf{AQ3 --- Controller contract and state representation}\par
\todo{Which admissible state sets, including heading and velocity limits, are
supported by the navigation primitives in each implementation? Confirm that each
primitive reaches its target's admissible set in finite time and stays within the
allowed regions. The formal contract is stated in Definition~2.2.}
\par\medskip
\noindent\textbf{AQ4 --- Relaxation example costs and source data}\par
\todo{Do the original numerical relaxation examples include the suffix closing
edge in both cost and violation? Confirm the underlying instances before
changing their reported numbers. The formulas and sufficient penalty bound
have been corrected in blue; the original example statistics are retained.}
\par\medskip
\noindent\textbf{AQ5 --- Assistance timing and optimality of detour evaluation}\par
\todo{In Section~3.2, are request times arrival times, action-start times, or
completion times, and does the confirmation time include all action
durations? Also, a shortest route to an assisting state need not minimize
the absolute timing-mismatch objective when arriving later reduces that
mismatch. Is waiting modeled explicitly in \texttt{EvalReq} and
\texttt{DelayCol}, and how should the timing and optimality assumptions be stated
for the candidate paths actually considered?}
\par\medskip
\noindent\textbf{AQ6 --- R-poset pruning and fleet-search guarantees}\par
\todo{What restrictions on formulas and automata justify sequentializing a
simultaneous transition, guarantee termination of repeated decompositions,
and ensure that the relaxed relations accept only valid words? Individual
swap tests alone need not establish arbitrary combinations of relaxations.
What state-equivalence condition permits dominance pruning by the scalar
distance/capacity profile in Section~4.2, and what bounds search depth in
automata containing cycles? The pruning, R-poset, and fleet-search
theorems are left unchanged until their source assumptions are fixed.}
\par\medskip
\noindent\textbf{AQ7 --- Communication-event refinement pseudocode}\par
\todo{In \texttt{ComOpt}, the two time variables are initialized equal, so the
displayed while condition is false at entry. The initial common time also
does not account for every robot's travel to the common location. What
initial feasible schedule and convergence variable does the implementation
use? Moving a graph articulation vertex while retaining one anchor link
does not alone preserve connectivity: is a spanning-tree or full-connectivity
test intended? The algorithm and connectivity lemma are left unchanged.}
\par\medskip
\noindent\textbf{AQ8 --- Latency invariant and infinite-horizon objective}\par
\todo{In Section~5.2, when the earliest possible return already exceeds the
latency deadline, how can appending that return restore the claimed bound?
Specify the prior feasibility invariant, return-reservation rule, and the
meaning of each time stamp. Separately, stationarity and ergodicity alone
do not make independent per-cycle rate maximization globally optimal when
the decisions change future states and cycle lengths. Which renewal or
separability assumptions support the final claim in Section~5.1, or should
it be weakened to a local approximation? These passages remain unchanged.}
\par\medskip
\noindent\textbf{AQ9 --- Probabilistic satisfaction and safe return}\par
\todo{In Section~6.1, summing the indicator of an absorbing accepting set
counts it infinitely often after entry; it does not count one hit. Is a
first-hit reward or stopped process intended in the task and return
equations? For a fixed arbitrary policy, merely reaching an AMEC does not
ensure that the policy follows accepting actions thereafter. The safe-return
expression also sums return values over all switching times rather than
specifying a switching law or a pointwise lower bound. Which notion is
intended? The return formula contains $\Diamond\Box$, so its later
description as sc-LTL also needs correction together with that choice.
These definitions, dependent optimizations, and guarantees are left
unchanged to avoid selecting a different safety problem without the authors.}
\par\medskip
\noindent\textbf{AQ10 --- Reactive task semantics and risk conventions}\par
\todo{Does Section~6.2 intend a general LTL recurring response rule with
finite sc-LTL obligations released online, and what certifies termination?
The current text calls the complete rule sc-LTL and uses $\varphi(t)$
both for a set and for a conjunction. In addition, the displayed VaR
definition uses an infimum of thresholds satisfying an upper-tail
probability lower bound, which is unbounded below. Is $\alpha$ a tail
probability or a confidence level, and should the standard quantile or
optimization definition of CVaR be used? These passages remain unchanged.}
\par\medskip
\noindent\textbf{AQ11 --- Prediction coverage and target-speed assumptions}\par
\todo{What independence or dependence assumptions justify multiplying the
per-subtask coverage probabilities in Lemma~6.3? Without independence,
a union-bound treatment is needed, with special care when several subtasks
share a target or time. The largest observed target speed is not a bound
on future speed without an additional model assumption. Finally, prediction
regions alone do not specify a sampling distribution for CVaR evaluation.
Which sampling law and speed bound are used? The coverage, delay, and
risk claims remain unchanged pending these choices.}
\par\medskip
\noindent\textbf{AQ12 --- Partial-team robustness in the simulations}\par
\todo{Confirm whether the simulations used endpoint-conditioned prior means
for unplanned robots when computing the local task cost. Confirm how any
additional connectivity margin was set in the implementation.}
\par\medskip
\noindent\textbf{AQ13 --- Experimental provenance and remaining proof assumptions}\par
\todo{Can the authors supply the underlying logs or scripts for the reproduced
and hardware statistics before claiming any new numerical verification?
Also confirm the additional assumptions needed for monotone coalition-cost
decrease in Section~4.2 when targets move, and the constraints used to
certify all collision and signal requirements in Section~6.3. The supplied
archive contains source, figures, and a compiled paper, but no experiment
code or raw logs; the reported results have therefore been preserved.}

\newpage\par\medskip
\noindent\textbf{BIBQ1 --- Unresolved unused reference records}\par
\todo{For each of the 74 records in
\texttt{REFERENCE\_AUTHOR\_QUESTIONS.md}, provide the exact DOI,
publisher or proceedings record, original document, or official repository.
The original record text is preserved unchanged outside the active
bibliography. Which work and version was intended?}

\par\medskip
\noindent\textbf{BIBQ2 --- Replaced citation records}\par
\todo{Confirm the seven verified replacements in the reference audit table.
If an original record was intended to identify another actual paper,
provide its authoritative record before restoring that citation.}

\par\medskip
\noindent\textbf{BIBQ3 --- Retracted unused paper}\par
\todo{The unused entry \texttt{abdel2018hybrid} is confirmed retracted.
It is preserved for provenance and excluded from scientific support.
Confirm that it should remain excluded from future versions.}

\par\medskip
\noindent\textbf{BIBQ4 --- Consulted editions and MoRoCo version}\par
\todo{Confirm the consulted editions for Latombe, Godsil and Royle,
and Sontag, whose original publication dates replace later distribution
years. Confirm that Chapter~5 is intended to follow MoRoCo arXiv version~1,
as its title and authors indicate, or specify the technical changes needed
before switching to version~2. No chapter has been silently changed to
follow a different version.}

\end{document}
